% !TeX program = xelatex
% UTF-8; compile twice with XeLaTeX. No external images or data files needed.
\documentclass[a4paper,landscape,10pt]{article}
\usepackage[margin=14mm,headheight=14pt,headsep=5mm,footskip=7mm]{geometry}
\usepackage{fontspec}
\setmainfont{DejaVu Sans}
\newfontfamily\ipafont{DejaVu Sans}
\usepackage[table]{xcolor}
\usepackage{array,longtable,booktabs,fancyhdr}
\definecolor{ink}{HTML}{18384D}
\definecolor{band}{HTML}{E7EFF4}
\pagestyle{fancy}\fancyhf{}
\fancyhead[L]{\small\color{ink}Parlez Vous QSO?! · Português / Portuguese}
\fancyhead[R]{\small Englische Phrase · Tradução portuguesa · IPA}
\fancyfoot[L]{\scriptsize Neufassung: europäisches Portugiesisch · Gliederung nach Spanish v 1.1}
\fancyfoot[R]{\small\thepage}
\setlength{\parindent}{0pt}
\setlength{\tabcolsep}{3mm}
\renewcommand{\arraystretch}{1.25}
\newcommand{\ipa}[1]{{\ipafont[#1]}}
\begin{document}
{\LARGE\bfseries\color{ink}Parlez Vous QSO?!}\hfill{\large Português / Portuguese}

\vspace{2mm}
{\small QSO-Sprachführer mit portugiesischer IPA-Lautschrift. Die Lautschrift bezieht sich auf die portugiesische Spalte.}

\vspace{2mm}
{\footnotesize \textbf{Leseschlüssel:} \ipa{ˈ} Hauptbetonung; \ipa{ʃ} wie „sch“, \ipa{ʒ} wie „j“ in „Journal“;
\ipa{ɲ} wie spanisches „ñ“, \ipa{ʎ} ein palataler l-Laut, \ipa{ɾ} ein kurzer r-Laut.
\ipa{ɨ} ist ein schwacher ungerundeter Vokal; \ipa{ɐ̃w̃} und \ipa{ɐ̃j̃} sind nasale Diphthonge. Grundlage: Portugal.}

\vspace{1mm}
\fontsize{9}{11.8}\selectfont
\begin{longtable}{@{}>{\raggedright\arraybackslash}p{13mm}>{\raggedright\arraybackslash}p{73mm}>{\raggedright\arraybackslash}p{75mm}>{\raggedright\arraybackslash}p{90mm}@{}}
\toprule
\textbf{Ref.} & \textbf{English phrase} & \textbf{Português} & \textbf{Lautschrift (IPA)}\\
\midrule
\endfirsthead
\toprule
\textbf{Ref.} & \textbf{English phrase} & \textbf{Português} & \textbf{Lautschrift (IPA)}\\
\midrule
\endhead
\midrule
\multicolumn{4}{r}{\scriptsize Fortsetzung auf der nächsten Seite}\\
\endfoot
\bottomrule
\endlastfoot
\rowcolor{band} 1.0 & \textbf{Opening Greetings} & \multicolumn{2}{l}{\textbf{Saudações iniciais}} \\*
1.1 & I am very pleased to meet you. & Tenho muito gosto em conhecer-te. & \ipa{ˈtɐɲu ˈmũj̃tu ˈɡoʃtu ɐ̃j̃ kuɲɨˈseɾ tɨ} \\
1.2 & Thank you for your call & Obrigado pela tua chamada. & \ipa{ubɾiˈɡadu ˈpɛlɐ ˈtuɐ ʃɐˈmadɐ} \\
1.3 & Good morning & Bom dia. & \ipa{bõ ˈdiɐ} \\
1.31 & Good day & Bom dia. & \ipa{bõ ˈdiɐ} \\
1.32 & Good afternoon & Boa tarde. & \ipa{ˈboɐ ˈtaɾdɨ} \\
1.33 & Good evening & Boa noite. & \ipa{ˈboɐ ˈnojtɨ} \\
1.4 & I know only a little Portuguese but I enjoy QSOs in Portuguese. & Sei só um pouco de português, mas gosto de fazer QSOs em português. & \ipa{sɐj sɔ ũ ˈpoku dɨ puɾtuˈɡeʃ maʃ ˈɡɔʃtu dɨ fɐˈzeɾ ke ˈɛsɨ ɔʃ ɐ̃j̃ puɾtuˈɡeʃ} \\
\rowcolor{band} 2.0 & \textbf{Operator} & \multicolumn{2}{l}{\textbf{Operador}} \\*
2.1 & My name is Peter. & Chamo-me Peter. & \ipa{ˈʃɐmu mɨ ˈpiːtə} \\
2.2 & I have been licensed since 1960. & Tenho licença desde 1960. & \ipa{ˈtɐɲu liˈsẽsɐ ˈdɛʒdɨ miɫ nuvɨˈsẽtuʃ i sɨˈsẽtɐ} \\
2.3 & I have been licensed for 43 years. & Tenho licença há 43 anos. & \ipa{ˈtɐɲu liˈsẽsɐ a kwɐˈɾẽtɐ i tɾeʃ ˈɐnuʃ} \\
2.4 & I live in Sittingbourne. & Moro em Sittingbourne. & \ipa{ˈmɔɾu ɐ̃j̃ ˈsɪtɪŋbɔːn} \\
2.5 & Sittingbourne is a village. & Sittingbourne é uma aldeia. & \ipa{ˈsɪtɪŋbɔːn ɛ ˈumɐ aɫˈdɐjɐ} \\
2.51 & Sittingbourne is a town. & Sittingbourne é uma vila. & \ipa{ˈsɪtɪŋbɔːn ɛ ˈumɐ ˈvilɐ} \\
2.52 & Sittingbourne is a city. & Sittingbourne é uma cidade. & \ipa{ˈsɪtɪŋbɔːn ɛ ˈumɐ siˈdadɨ} \\
2.6 & Sittingbourne is between London and Dover. & Sittingbourne fica entre Londres e Dover. & \ipa{ˈsɪtɪŋbɔːn ˈfikɐ ˈẽtɾɨ ˈlõdɾɨʃ i ˈdəʊvə} \\
2.7 & I live in the South of England. & Moro no sul de Inglaterra. & \ipa{ˈmɔɾu nu suɫ dɨ ĩɡlɐˈtɛʁɐ} \\
2.71 & I live in the North of England. & Moro no norte de Inglaterra. & \ipa{ˈmɔɾu nu ˈnɔɾtɨ dɨ ĩɡlɐˈtɛʁɐ} \\
2.72 & I live in the East of England. & Moro no leste de Inglaterra. & \ipa{ˈmɔɾu nu ˈlɛʃtɨ dɨ ĩɡlɐˈtɛʁɐ} \\
2.73 & I live in West of England. & Moro no oeste de Inglaterra. & \ipa{ˈmɔɾu nu oˈɛʃtɨ dɨ ĩɡlɐˈtɛʁɐ} \\
2.74 & I live in South East England. & Moro no sudeste de Inglaterra. & \ipa{ˈmɔɾu nu suˈdɛʃtɨ dɨ ĩɡlɐˈtɛʁɐ} \\
2.75 & I live in South West England. & Moro no sudoeste de Inglaterra. & \ipa{ˈmɔɾu nu sudoˈɛʃtɨ dɨ ĩɡlɐˈtɛʁɐ} \\
2.76 & I live in North East England. & Moro no nordeste de Inglaterra. & \ipa{ˈmɔɾu nu nuɾˈdɛʃtɨ dɨ ĩɡlɐˈtɛʁɐ} \\
2.77 & I live in North West England. & Moro no noroeste de Inglaterra. & \ipa{ˈmɔɾu nu nuɾoˈɛʃtɨ dɨ ĩɡlɐˈtɛʁɐ} \\
2.78 & I live in central England. & Moro no centro de Inglaterra. & \ipa{ˈmɔɾu nu ˈsẽtɾu dɨ ĩɡlɐˈtɛʁɐ} \\
2.8 & I live in the country. & Moro no campo. & \ipa{ˈmɔɾu nu ˈkɐ̃pu} \\
2.81 & I live by the sea. & Moro junto ao mar. & \ipa{ˈmɔɾu ˈʒũtu aw maɾ} \\
2.82 & I live near mountains. & Moro perto das montanhas. & \ipa{ˈmɔɾu ˈpɛɾtu dɐʃ mõˈtɐɲɐʃ} \\
2.9 & Sittingbourne is 50Km from London. & Sittingbourne fica a 50 km de Londres. & \ipa{ˈsɪtɪŋbɔːn ˈfikɐ ɐ sĩˈkwẽtɐ kiˈlɔmɨtɾuʃ dɨ ˈlõdɾɨʃ} \\
2.11 & I am retired. & Estou reformado. & \ipa{ɨʃˈtow ʁɨfuɾˈmadu} \\
2.12 & My home call sign is M0ECS & O meu indicativo de origem é M0ECS. & \ipa{u mew ĩdikɐˈtivu dɨ uˈɾiʒɐ̃j̃ ɛ ˈɛmɨ ˈzɛɾu ɛ se ˈɛsɨ} \\
\rowcolor{band} 3.0 & \textbf{Station} & \multicolumn{2}{l}{\textbf{Estação}} \\*
3.1 & My radio is an Icom 706. & O meu rádio é um Icom 706. & \ipa{u mew ˈʁadju ɛ ũ ˈaɪkɒm sɛtɨˈsẽtuʃ i sɐjʃ} \\
3.2 & My station is all homebrew (homemade). & A minha estação é toda de construção caseira. & \ipa{ɐ ˈmiɲɐ ɨʃtɐˈsɐ̃w̃ ɛ ˈtodɐ dɨ kõʃtɾuˈsɐ̃w̃ kɐˈzɐjɾɐ} \\
3.3 & I am running 5 Watts. & Estou a transmitir com 5 watts. & \ipa{ɨʃˈtow ɐ tɾɐ̃ʒmiˈtiɾ kõ ˈsĩku watʃ} \\
3.4 & My antenna is a long wire. & A minha antena é de fio longo. & \ipa{ɐ ˈmiɲɐ ɐ̃ˈtenɐ ɛ dɨ ˈfiu ˈlõɡu} \\
3.41 & My antenna is a dipole. & A minha antena é um dipolo. & \ipa{ɐ ˈmiɲɐ ɐ̃ˈtenɐ ɛ ũ diˈpɔlu} \\
3.42 & My antenna is a G5RV. & A minha antena é uma G5RV. & \ipa{ɐ ˈmiɲɐ ɐ̃ˈtenɐ ɛ ˈumɐ ʒe ˈsĩku ˈɛʁɨ ve} \\
3.43 & My antenna is a doublet. & A minha antena é um dipolo do tipo doublet. & \ipa{ɐ ˈmiɲɐ ɐ̃ˈtenɐ ɛ ũ diˈpɔlu du ˈtipu ˈdʌblɪt} \\
3.44 & My antenna is a delta loop. & A minha antena é uma delta loop. & \ipa{ɐ ˈmiɲɐ ɐ̃ˈtenɐ ɛ ˈumɐ ˈdɛɫtɐ luːp} \\
3.45 & My antenna is a vertical. & A minha antena é vertical. & \ipa{ɐ ˈmiɲɐ ɐ̃ˈtenɐ ɛ vɨɾtiˈkaɫ} \\
3.46 & My antenna is a log periodic. & A minha antena é log-periódica. & \ipa{ɐ ˈmiɲɐ ɐ̃ˈtenɐ ɛ lɔɡ pɨˈɾjɔdikɐ} \\
3.47 & My antenna is a Yagi beam. & A minha antena é uma Yagi direcional. & \ipa{ɐ ˈmiɲɐ ɐ̃ˈtenɐ ɛ ˈumɐ ˈjaɡi diɾɨsjuˈnaɫ} \\
3.48 & My antenna is a 5 element beam. & A minha antena é direcional, com 5 elementos. & \ipa{ɐ ˈmiɲɐ ɐ̃ˈtenɐ ɛ diɾɨsjuˈnaɫ kõ ˈsĩku ilɨˈmẽtuʃ} \\
\rowcolor{band} 4.0 & \textbf{Conditions} & \multicolumn{2}{l}{\textbf{Condições de propagação}} \\*
4.1 & Thank you for my 599 report. & Obrigado pelo controlo de sinal: 599. & \ipa{ubɾiˈɡadu ˈpɛlu kõˈtɾolu dɨ siˈnaɫ ˈsĩku ˈnɔvɨ ˈnɔvɨ} \\
4.2 & The band conditions are good. & As condições de propagação na banda são boas. & \ipa{ɐʃ kõdiˈsõj̃ʃ dɨ pɾupɐɡɐˈsɐ̃w̃ nɐ ˈbɐ̃dɐ sɐ̃w̃ ˈboɐʃ} \\
4.3 & The band conditions are bad. & As condições de propagação na banda são más. & \ipa{ɐʃ kõdiˈsõj̃ʃ dɨ pɾupɐɡɐˈsɐ̃w̃ nɐ ˈbɐ̃dɐ sɐ̃w̃ maʃ} \\
4.4 & There is strong interference here. & Há muita interferência aqui. & \ipa{a ˈmũj̃tɐ ĩtɨɾfɨˈɾẽsjɐ ɐˈki} \\
4.5 & Can you change to the 80m band? & Podes mudar para a banda dos 80 metros? & \ipa{ˈpɔdɨʃ muˈdaɾ ˈpɐɾɐ ɐ ˈbɐ̃dɐ duʃ ojˈtẽtɐ ˈmɛtɾuʃ} \\
4.6 & Please change frequency to 3.558 MHz. & Por favor, muda para a frequência de 3,558 MHz. & \ipa{puɾ fɐˈvoɾ ˈmudɐ ˈpɐɾɐ ɐ fɾɨˈkwẽsjɐ dɨ tɾeʃ ˈviɾɡulɐ ˈsĩku ˈsĩku ˈojtu ˈmɛɡɐ ɛɾtʃ} \\
\rowcolor{band} 5.0 & \textbf{General Queries and Comments} & \multicolumn{2}{l}{\textbf{Perguntas e comentários gerais}} \\*
5.1 & What does <\#\#\#\#\#> mean please? & O que significa <\#\#\#\#\#>, por favor? & \ipa{u kɨ siɡniˈfikɐ … puɾ fɐˈvoɾ} \\
5.2 & Please wait a moment, I need to check my Portuguese dictionary. & Espera um momento, por favor. Preciso de consultar o meu dicionário de português. & \ipa{ɨʃˈpɛɾɐ ũ muˈmẽtu puɾ fɐˈvoɾ pɾɨˈsizu dɨ kõsuɫˈtaɾ u mew disjuˈnaɾju dɨ puɾtuˈɡeʃ} \\
5.3 & How do you say, <\#\#\#\#\#> in Portuguese ? & Como se diz <\#\#\#\#\#> em português? & \ipa{ˈkomu sɨ diʃ … ɐ̃j̃ puɾtuˈɡeʃ} \\
5.5 & My favourite modes of operation are CW / RTTY / PSK / MT63 / Olivia / Hellschreiber / Digital Modes / and SSB & Os meus modos preferidos são CW / RTTY / PSK / MT63 / Olivia / Hellschreiber / modos digitais / e SSB. & \ipa{uʃ mewʃ ˈmɔduʃ pɾɨfɨˈɾiduʃ sɐ̃w̃ se ˈdɐbɫju / ˈɛʁɨ te te ˈipsɨlɔn / pe ˈɛsɨ ka / ˈɛmɨ te sɨˈsẽtɐ i tɾeʃ / oˈlivjɐ / ˈhɛlˌʃʁaɪbɐ / ˈmɔduʃ diʒiˈtajʃ / i ˈɛsɨ ˈɛsɨ be} \\
5.6 & My best DX is Australia. & O meu melhor DX é a Austrália. & \ipa{u mew mɨˈʎɔɾ de ʃiʃ ɛ ɐ awʃˈtɾaljɐ} \\
5.7 & I enjoy working contests. & Gosto de participar em concursos. & \ipa{ˈɡɔʃtu dɨ pɐɾtisiˈpaɾ ɐ̃j̃ kõˈkuɾsuʃ} \\
5.8 & I like working QRP. & Gosto de trabalhar em QRP. & \ipa{ˈɡɔʃtu dɨ tɾɐbɐˈʎaɾ ɐ̃j̃ ke ˈɛʁɨ pe} \\
\rowcolor{band} 6.0 & \textbf{The Weather} & \multicolumn{2}{l}{\textbf{O tempo}} \\*
6.1 & It is sunny. & Está sol. & \ipa{ɨʃˈta sɔɫ} \\
6.2 & It is windy. & Está vento. & \ipa{ɨʃˈta ˈvẽtu} \\
6.3 & It is raining. & Está a chover. & \ipa{ɨʃˈta ɐ ʃuˈveɾ} \\
6.4 & It is snowing. & Está a nevar. & \ipa{ɨʃˈta ɐ nɨˈvaɾ} \\
6.5 & It is very hot. & Está muito calor. & \ipa{ɨʃˈta ˈmũj̃tu kɐˈloɾ} \\
6.6 & It is very cold. & Está muito frio. & \ipa{ɨʃˈta ˈmũj̃tu ˈfɾiu} \\
6.7 & The temperature is 8C. & A temperatura é de 8 °C. & \ipa{ɐ tɨ̃pɨɾɐˈtuɾɐ ɛ dɨ ˈojtu ˈɡɾawʃ ˈsɛɫsjuʃ} \\
6.8 & It is overcast. & O céu está encoberto. & \ipa{u sɛw ɨʃˈta ẽkuˈbɛɾtu} \\
6.9 & It is foggy. & Há nevoeiro. & \ipa{a nɨvuˈɐjɾu} \\
6.11 & It has been raining all day. & Tem estado a chover o dia todo. & \ipa{tɐ̃j̃ ɨʃˈtadu ɐ ʃuˈveɾ u ˈdiɐ ˈtodu} \\
6.12 & There is 20cm of snow here. & Há 20 cm de neve aqui. & \ipa{a ˈvĩtɨ sẽˈtimɨtɾuʃ dɨ ˈnɛvɨ ɐˈki} \\
\rowcolor{band} 8.0 & \textbf{QSL} & \multicolumn{2}{l}{\textbf{Cartões QSL}} \\*
8.1 & My QSL card is sure via the bureau & O meu cartão QSL está garantido via bureau. & \ipa{u mew kɐɾˈtɐ̃w̃ ke ˈɛsɨ ˈɛlɨ ɨʃˈta ɡɐɾɐ̃ˈtidu ˈviɐ byˈʁo} \\
8.2 & QSL via eQSL & QSL via eQSL. & \ipa{ke ˈɛsɨ ˈɛlɨ ˈviɐ i ke ˈɛsɨ ˈɛlɨ} \\
8.3 & No eQSL please & Por favor, não envies eQSL. & \ipa{puɾ fɐˈvoɾ nɐ̃w̃ ẽˈviɨʃ i ke ˈɛsɨ ˈɛlɨ} \\
8.4 & I will send you my card direct. & Vou enviar-te o meu cartão diretamente. & \ipa{vow ẽviˈaɾ tɨ u mew kɐɾˈtɐ̃w̃ diɾɛtɐˈmẽtɨ} \\
8.5 & What is your address? & Qual é a tua morada? & \ipa{kwaɫ ɛ ɐ ˈtuɐ muˈɾadɐ} \\
8.6 & Please send your QSL card direct. & Por favor, envia-me o teu cartão QSL diretamente. & \ipa{puɾ fɐˈvoɾ ẽˈviɐ mɨ u tew kɐɾˈtɐ̃w̃ ke ˈɛsɨ ˈɛlɨ diɾɛtɐˈmẽtɨ} \\
8.7 & My address is... & A minha morada é... & \ipa{ɐ ˈmiɲɐ muˈɾadɐ ɛ …} \\
8.8 & My address is shown on QSL.com / EHam.net & A minha morada está em QSL.com / EHam.net. & \ipa{ɐ ˈmiɲɐ muˈɾadɐ ɨʃˈta ɐ̃j̃ ke ˈɛsɨ ˈɛlɨ ˈpõtu kɔm / i hæm ˈpõtu nɛt} \\
8.9 & Sorry I do not send QSL cards. & Lamento, mas não envio cartões QSL. & \ipa{lɐˈmẽtu maʃ nɐ̃w̃ ẽˈviu kɐɾˈtõj̃ʃ ke ˈɛsɨ ˈɛlɨ} \\
8.11 & Please send QSL card via G1JQH. & Por favor, envia o cartão QSL via G1JQH. & \ipa{puɾ fɐˈvoɾ ẽˈviɐ u kɐɾˈtɐ̃w̃ ke ˈɛsɨ ˈɛlɨ ˈviɐ ʒe ũ ˈʒɔtɐ ke ɐˈɡa} \\
8.12 & My WACRAL / GQRP / DXCC / FISTS number is 1234. & O meu número WACRAL / GQRP / DXCC / FISTS é 1234. & \ipa{u mew ˈnumɨɾu ˈdɐbɫju a se ˈɛʁɨ a ˈɛlɨ / ʒe ke ˈɛʁɨ pe / de ʃiʃ se se / fɪsts ɛ miɫ duˈzẽtuʃ i ˈtɾĩtɐ i ˈkwatɾu} \\
\rowcolor{band} 9.0 & \textbf{Final Greetings} & \multicolumn{2}{l}{\textbf{Saudações finais}} \\*
9.1 & Thank you. & Obrigado. & \ipa{ubɾiˈɡadu} \\
9.2 & Thank you for a very enjoyable contact. & Obrigado por um contacto muito agradável. & \ipa{ubɾiˈɡadu puɾ ũ kõˈtaktu ˈmũj̃tu ɐɡɾɐˈdavɛɫ} \\
9.3 & I enjoyed our chat. & Gostei da nossa conversa. & \ipa{ɡuʃˈtɐj dɐ ˈnɔsɐ kõˈvɛɾsɐ} \\
9.4 & I enjoyed our QSO in Portuguese. & Gostei do nosso QSO em português. & \ipa{ɡuʃˈtɐj du ˈnɔsu ke ˈɛsɨ ɔ ɐ̃j̃ puɾtuˈɡeʃ} \\
9.5 & It has been great fun. & Foi muito divertido. & \ipa{foj ˈmũj̃tu divɨɾˈtidu} \\
9.6 & If you would like a regular contact with me, you can contact me via email. & Se quiseres manter um contacto regular comigo, podes contactar-me por correio eletrónico. & \ipa{sɨ kiˈzɛɾɨʃ mɐ̃ˈteɾ ũ kõˈtaktu ʁɨɡuˈlaɾ kuˈmiɡu ˈpɔdɨʃ kõtakˈtaɾ mɨ puɾ kuˈʁɐju ilɛˈtɾɔniku} \\
9.7 & My email address is... & O meu endereço de correio eletrónico é... & \ipa{u mew ẽdɨˈɾesu dɨ kuˈʁɐju ilɛˈtɾɔniku ɛ …} \\
9.8 & My email address is shown on QSL.com / EHam.net & O meu endereço de correio eletrónico está em QSL.com / EHam.net. & \ipa{u mew ẽdɨˈɾesu dɨ kuˈʁɐju ilɛˈtɾɔniku ɨʃˈta ɐ̃j̃ ke ˈɛsɨ ˈɛlɨ ˈpõtu kɔm / i hæm ˈpõtu nɛt} \\
9.9 & I look forward to meeting you again. & Espero voltar a encontrar-te. & \ipa{ɨʃˈpɛɾu voɫˈtaɾ ɐ ẽkõˈtɾaɾ tɨ} \\
9.11 & My best wishes to you and your family. & Os meus melhores votos para ti e para a tua família. & \ipa{uʃ mewʃ mɨˈʎɔɾɨʃ ˈvɔtuʃ ˈpɐɾɐ ti i ˈpɐɾɐ ɐ ˈtuɐ fɐˈmiljɐ} \\
9.12 & Best wishes to your wife. & Os meus melhores votos para a tua esposa. & \ipa{uʃ mewʃ mɨˈʎɔɾɨʃ ˈvɔtuʃ ˈpɐɾɐ ɐ ˈtuɐ ɨʃˈpozɐ} \\
9.121 & Best wishes to your husband. & Os meus melhores votos para o teu marido. & \ipa{uʃ mewʃ mɨˈʎɔɾɨʃ ˈvɔtuʃ ˈpɐɾɐ u tew mɐˈɾidu} \\
9.13 & Good luck. & Boa sorte. & \ipa{ˈboɐ ˈsɔɾtɨ} \\
9.14 & Good health. & Muita saúde. & \ipa{ˈmũj̃tɐ sɐˈudɨ} \\
9.15 & Goodbye & Adeus. & \ipa{ɐˈdewʃ} \\
9.16 & I wish you an enjoyable afternoon. & Desejo-te uma boa tarde. & \ipa{dɨˈzeʒu tɨ ˈumɐ ˈboɐ ˈtaɾdɨ} \\
9.161 & I wish you an enjoyable evening. & Desejo-te uma boa noite. & \ipa{dɨˈzeʒu tɨ ˈumɐ ˈboɐ ˈnojtɨ} \\
9.162 & I wish you an enjoyable weekend. & Desejo-te um bom fim de semana. & \ipa{dɨˈzeʒu tɨ ũ bõ fĩ dɨ sɨˈmɐnɐ} \\
9.17 & Enjoy your holiday. & Boas férias. & \ipa{ˈboɐʃ ˈfɛɾjɐʃ} \\
9.18 & Happy New Year. & Feliz Ano Novo. & \ipa{fɨˈliʃ ˈɐnu ˈnovu} \\
\rowcolor{band} 10.0 & \textbf{Country} & \multicolumn{2}{l}{\textbf{Países e territórios}} \\*
10.1 & Albania & Albânia & \ipa{aɫˈbɐnjɐ} \\
10.2 & Austria & Áustria & \ipa{ˈawʃtɾjɐ} \\
10.3 & Belgium & Bélgica & \ipa{ˈbɛɫʒikɐ} \\
10.4 & Belorussia & Bielorrússia & \ipa{bjɛluˈʁusjɐ} \\
10.5 & Bosnia \& Herzegovina & Bósnia e Herzegovina & \ipa{ˈbɔznjɐ i ɛɾzɨɡuˈvinɐ} \\
10.6 & Bulgaria & Bulgária & \ipa{buɫˈɡaɾjɐ} \\
10.7 & Croatia & Croácia & \ipa{kɾuˈasjɐ} \\
10.8 & Czech Republic & República Checa & \ipa{ʁɨˈpublikɐ ˈʃɛkɐ} \\
10.9 & Denmark & Dinamarca & \ipa{dinɐˈmaɾkɐ} \\
10.11 & Eire & Irlanda & \ipa{iɾˈlɐ̃dɐ} \\
10.12 & England & Inglaterra & \ipa{ĩɡlɐˈtɛʁɐ} \\
10.13 & Estonia & Estónia & \ipa{ɨʃˈtɔnjɐ} \\
10.14 & Finland & Finlândia & \ipa{fĩˈlɐ̃djɐ} \\
10.15 & France & França & \ipa{ˈfɾɐ̃sɐ} \\
10.16 & Germany & Alemanha & \ipa{ɐlɨˈmɐɲɐ} \\
10.17 & Gibraltar & Gibraltar & \ipa{ʒibɾaɫˈtaɾ} \\
10.18 & Greece & Grécia & \ipa{ˈɡɾɛsjɐ} \\
10.19 & Holland & Holanda & \ipa{uˈlɐ̃dɐ} \\
10.21 & Hungary & Hungria & \ipa{ũˈɡɾiɐ} \\
10.22 & Iceland & Islândia & \ipa{iʒˈlɐ̃djɐ} \\
10.23 & Northern Ireland & Irlanda do Norte & \ipa{iɾˈlɐ̃dɐ du ˈnɔɾtɨ} \\
10.24 & Isle of Man & Ilha de Man & \ipa{ˈiʎɐ dɨ mæn} \\
10.25 & Italy & Itália & \ipa{iˈtaljɐ} \\
10.26 & Latvia & Letónia & \ipa{lɨˈtɔnjɐ} \\
10.27 & Liechtenstein & Liechtenstein & \ipa{ˈliçtn̩ʃtaɪn} \\
10.28 & Lithuania & Lituânia & \ipa{liˈtwɐnjɐ} \\
10.29 & Luxembourg & Luxemburgo & \ipa{luksɐ̃ˈbuɾɡu} \\
10.31 & Macedonia & Macedónia do Norte & \ipa{mɐsɨˈdɔnjɐ du ˈnɔɾtɨ} \\
10.32 & Moldavia & Moldávia & \ipa{moɫˈdavjɐ} \\
10.33 & Netherlands & Países Baixos & \ipa{pɐˈizɨʃ ˈbajʃuʃ} \\
10.34 & Norway & Noruega & \ipa{nuɾuˈɛɡɐ} \\
10.35 & Poland & Polónia & \ipa{puˈlɔnjɐ} \\
10.36 & Portugal & Portugal & \ipa{puɾtuˈɡaɫ} \\
10.37 & Romania & Roménia & \ipa{ʁuˈmɛnjɐ} \\
10.38 & Russian Federation & Federação Russa & \ipa{fɨdɨɾɐˈsɐ̃w̃ ˈʁusɐ} \\
10.39 & Scotland & Escócia & \ipa{ɨʃˈkɔsjɐ} \\
10.41 & Slovakia & Eslováquia & \ipa{ɨʒluˈvakjɐ} \\
10.42 & Slovenia & Eslovénia & \ipa{ɨʒluˈvɛnjɐ} \\
10.43 & Spain & Espanha & \ipa{ɨʃˈpɐɲɐ} \\
10.44 & Svalbard & Svalbard & \ipa{ˈsvɑːɫbɑr} \\
10.45 & Sweden & Suécia & \ipa{suˈɛsjɐ} \\
10.46 & Switzerland & Suíça & \ipa{suˈisɐ} \\
10.47 & Turkey & Turquia & \ipa{tuɾˈkiɐ} \\
10.48 & Ukraine & Ucrânia & \ipa{uˈkɾɐnjɐ} \\
10.49 & Wales & País de Gales & \ipa{pɐˈiʃ dɨ ˈɡalɨʃ} \\
\end{longtable}
\clearpage
\normalsize
\section*{Hinweise zur Neufassung und Aussprache}
\fontsize{8.5}{10.5}\selectfont\raggedright
\textbf{Vorlage und Umfang.} Diese portugiesische Fassung ist neu entworfen; eine portugiesische
Originaldatei lag nicht vor. Als Strukturvorlage dient die bereits übertragene Fassung von
„Spanish v 1.1.doc“ mit den Credits Vicente -- EA5OO und Colin Bass (ohne Rufzeichen).
Diese Credits beziehen sich auf die spanische Vorlage, nicht auf die neue portugiesische Übersetzung.
Alle 142 Phrasen und Ländernamen sowie die ursprünglichen Referenzen und Nummernsprünge sind enthalten.
Der fehlende Abschnitt 7 wird nicht neu erfunden. Die englischen sprachbezogenen Phrasen nennen nun
„Portuguese“ statt „Spanish“. Namen, Rufzeichen, Modellnummern und Beispielwerte können durch eigene Angaben ersetzt werden.

\medskip
\textbf{Sprachvariante und Anrede.} Übersetzung und Lautschrift orientieren sich am europäischen
Portugiesisch, besonders am Raum Lissabon. Verwendet wird durchgehend die vertraute Anrede „tu“;
die Verben und die Formen „te“, „tua“ und „teu“ passen dazu. „Morada“, „reformado“ und
„correio eletrónico“ sind in Portugal übliche Formulierungen. „Obrigado“ und „reformado“ gelten
hier für einen männlichen Sprecher; eine Sprecherin verwendet „obrigada“ und „reformada“.
Die reguläre Schreibung nach dem portugiesischen Orthografieabkommen steht neben in Portugal
ausgesprochenen Konsonanten, etwa „contacto“ mit hörbarem \ipa{k}.

\medskip
\textbf{Bewusste Anpassungen.} In 8.12 wird einheitlich die Mitgliedsnummer 1234 verwendet, wie
in der englischen Spalte; die spanische Vorlage enthält dort abweichend 1224. Die Frequenz 3.558 MHz
erscheint portugiesisch als 3,558 MHz mit Dezimalkomma; es ist derselbe Wert.
„Boa tarde“ begrüßt am Nachmittag, „Boa noite“ am Abend und in der Nacht. Die starren Uhrzeitangaben
der spanischen Vorlage entfallen; auch die Abschiedswünsche folgen dieser Unterscheidung.
„Bom dia“ deckt sowohl „Good morning“ als auch das allgemeine „Good day“ ab.
In 6.11 beschreibt „Tem estado a chover o dia todo“ den bisherigen Regenverlauf.
Die Einträge „aldeia“, „vila“ und „cidade“ bleiben Alternativen zum Einsetzen des eigenen Orts;
sie sind keine drei gleichzeitigen Aussagen über Sittingbourne.

\medskip
\textbf{Amateurfunkbegriffe.} „Indicativo de origem“ (2.12) bezeichnet das Heimatrufzeichen,
„controlo de sinal“ (4.1) den Signalrapport. „De construção caseira“ steht für Eigenbau.
Bei Antennen werden „fio longo“, „dipolo“, „vertical“, „log-periódica“ und „direcional“ benutzt.
„Doublet“ bleibt als Typbezeichnung eines Dipols erhalten; „delta loop“ bezeichnet die dreieckige
Schleifenantenne. G5RV und Yagi bleiben Modell- beziehungsweise Bauartnamen.
„Via bureau“ meint den QSL-Versand über den Vermittlungsdienst des Amateurfunkverbands;
„diretamente“ meint den direkten Versand an den Empfänger. „Trabalhar em QRP“ bezeichnet
Funkbetrieb mit geringer Leistung. Die Auswahl der üblichen Fachbegriffe wurde mit portugiesischen
Amateurfunktexten der REP abgeglichen.

\medskip
\textbf{Ländernamen.} Die englischen Quellbezeichnungen bleiben nachvollziehbar erhalten.
Portugiesisch steht in 10.31 der heutige Staatsname „Macedónia do Norte“ für „Macedonia“.
„Bielorrússia“ und „Moldávia“ sind gebräuchliche portugiesische Namen.
„Holanda“ bleibt als Entsprechung für „Holland“ stehen; das gesamte Land heißt „Países Baixos“ (10.33).
Die Liste enthält neben Staaten auch Gebiete und Landesteile; daher lautet die Überschrift
„Países e territórios“. „Eire“ wird als „Irlanda“ wiedergegeben.

\medskip
\textbf{Transkriptionskonvention.} Die IPA-Spalte wurde manuell als breite Aussprachehilfe erstellt,
nicht als enge Transkription einer Aufnahme. Die Wörter werden weitgehend einzeln transkribiert;
Lautverbindungen über Wortgrenzen, Vokalausfall, Satzmelodie und Entstimmung werden nicht abgebildet.
Insbesondere steht auslautendes s/z grundsätzlich als \ipa{ʃ}, auch wenn es im flüssigen Satz
vor stimmhaften Konsonanten \ipa{ʒ} und vor Vokalen \ipa{z} werden kann.
\ipa{ɐ} bezeichnet einen zentralen a-Laut, \ipa{ɨ} den häufig stark reduzierten unbetonten e-Laut.
Die Tilde markiert Nasalierung: \ipa{ɐ̃, ẽ, ĩ, õ, ũ}; auch beide Bestandteile nasaler Diphthonge
sind markiert. \ipa{ɫ} steht für dunkles l am Silbenende; vor Vokalen wird vereinfachend \ipa{l} geschrieben.
\ipa{ʁ} ist ein hinterer r-Laut, \ipa{ɾ} ein kurzer Zungenschlag.
Kontextabhängige weichere Varianten von b, d und g werden nicht gesondert notiert.

\medskip
\textbf{Kürzel, Eigennamen und Zahlen.} Kürzel und Rufzeichen werden überwiegend mit portugiesischen
Buchstabennamen gelesen: Q „quê“, W „dâblio“, Y „ípsilon“, X „xis“.
Das ist keine Wiedergabe des internationalen Buchstabieralphabets.
FISTS bleibt ein englisch gesprochenes Wort, ebenso Peter, Sittingbourne, Dover, Icom, doublet und loop;
Hellschreiber und Liechtenstein sind deutsch wiedergegeben, Svalbard mit einer norwegischen Aussprachehilfe.
Das e in eQSL und EHam wird als englisches \ipa{i} gelesen. „Bureau“ ist französisch transkribiert.
599 wird ziffernweise „cinco nove nove“ gelesen, 3,558 „três vírgula cinco cinco oito“;
1234 „mil duzentos e trinta e quatro“, 706 „setecentos e seis“ und 1960 „mil novecentos e sessenta“.
„Watts“ und „megahertz“ sind hier als im Portugiesischen verwendete Lehnwörter gesprochen.
Schrägstriche trennen Alternativen oder Listenelemente; \texttt{<\#\#\#\#\#>} ist ein Platzhalter
und erscheint in der Lautschrift als Auslassung. Die Punkte in Webadressen heißen „ponto“.

\medskip
\textbf{Brasilien und Quellen.} Die Sätze sind meist auch in Brasilien verständlich; dort sind beispielsweise
„endereço“ statt „morada“, „aposentado“ statt „reformado“ und „e-mail“ gebräuchlicher.
Anrede und Aussprache unterscheiden sich regional deutlich. Die IPA-Spalte beschreibt Portugal;
einzelne Zeichen zu ersetzen ergibt keine vollständige brasilianische Fassung.
Phonetische Grundlage: Madalena Cruz-Ferreira, „European Portuguese“,
Journal of the International Phonetic Association 25(2), 1995, S. 90--94,
DOI: 10.1017/S0025100300005223. Fachsprachlicher Abgleich: REP / CT3KN,
„informação diversa“, Abschnitt Radioamadorismo, www.rep.pt/ct3kn/radioamadorismo.htm.
Einzelne Wortaussprachen wurden mit dem Dicionário Infopédia da Língua Portuguesa (Porto Editora),
www.infopedia.pt/dicionarios/lingua-portuguesa, abgeglichen.
Diese Quellen erläutern Aussprache und Fachsprache; sie liefern keine Übersetzung dieses Sprachführers.

\medskip
\textbf{LaTeX.} Selbstständige UTF-8-Datei, zweimal mit XeLaTeX kompilieren.
Benötigt werden DejaVu Sans und die üblichen LaTeX-Pakete der Präambel.
Die IPA-Zeichen sind direkt als Unicode eingegeben; externe Dateien sind nicht erforderlich.
\end{document}
